\sheet{3}{OpenMP scheduling, tasks, and synchronization}{1h30}
\section{Preparation}
Define \textbf{barrier}, \textbf{implicit barrier}, \textbf{task},
\textbf{taskwait}, \textbf{cutoff}, \textbf{false sharing} and \textbf{data race}.
For a barrier and a taskwait, identify which work must finish before execution
continues. Give an example of correct shared-memory code with poor performance.

Sources are in \code{03-openmp-advanced}. Each section has a separate
program and can be built before the others are complete. The supplied main
checks numerical results before measuring. Run commands from this \code{03-openmp-advanced} directory:
\begin{lstlisting}[style=shell]
cmake -S . -B build -DCMAKE_BUILD_TYPE=Release
\end{lstlisting}
Use a team size within the CPU allocation and rebuild the relevant target after
editing. Complete the TODO functions in the named file; retain its supplied
input generation and checks.

This folder contains its own CMake files and utility headers; it can be copied
and compiled independently. Numerical sizes are freely chosen integers within
the ranges listed in \code{README.md}; the examples are not a predefined menu.
For programs whose usage lists \code{--self-test}, add that flag to run extra
small and edge cases before the requested input. Without it, only the requested
size is checked and measured. Rebuild after changing the implementation.

\section{Exercises}
\exercise{C1}{Scheduling uneven work}
Edit \code{C1-schedule.cpp}. The supplied \code{irregular_value} performs 4,000
unsigned arithmetic updates for indices in the first quarter, and 32 for the
other indices. Equal iteration counts therefore need not mean equal work.
Unsigned wraparound is intentional; leave the recurrence unchanged.
\begin{enumerate}
\item On paper, schedule costs $(4,1,1,4,1,1)$ on two threads using
\code{static,1} and \code{dynamic,1}. For dynamic scheduling, a free thread takes
the next iteration; state a tie-breaking rule. Compare completion times without
assignment overhead. Is dynamic scheduling necessarily faster in this example?
\item Implement \code{scheduled_work}. Distribute indices with
\code{schedule(runtime)} and write \code{irregular_value(i,n)} to \code{out[i]}.
Each output must have one writer.
\item Build and run:
\begin{lstlisting}[style=shell]
cmake --build build --target c1_schedule -j
OMP_NUM_THREADS=4 OMP_DYNAMIC=false OMP_SCHEDULE=static \
  ./build/bin/c1_schedule 16384 5
\end{lstlisting}
The arguments are length and trials. Repeat with \code{OMP_SCHEDULE} set to
\code{static,1}, \code{dynamic,1} and \code{guided,8}. Keep input and placement fixed.
\item Compare medians and ranges. Explain the tradeoff between imbalance,
assignment overhead and locality. The checks compare every output, because
changing its assigned thread must not change the required value.
\end{enumerate}

\exercise{C2}{Private histograms and merging}
Edit \code{C2-histogram.cpp}. Input \code{x[i]} contributes to bin
\lstinline|x[i] % bins|. Different indices can select the same bin, so an
unsynchronized shared increment is unsafe. Separate per-thread histograms allow
independent counting followed by a merge.
\begin{enumerate}
\item Compute the histogram of \code{{0,1,7,2,8}} with three bins.
Explain why the total can be correct even if individual bins are wrong.
\item Implement \code{histogram}. Allocate zeroed storage for participating
threads, distribute input indices, and let each thread update its own row.
Finish all counting before merging rows into the returned vector.
\item Return exactly \code{bins} counters, including zeros. Empty input must
produce an all-zero vector. If allocating for a maximum team size, do not merge
uninitialised unused rows. Explain which barrier makes merging safe.
\item Build and run:
\begin{lstlisting}[style=shell]
cmake --build build --target c2_histogram -j
OMP_NUM_THREADS=4 OMP_DYNAMIC=false ./build/bin/c2_histogram 1048576 31 5
\end{lstlisting}
The arguments are length, bin count and trials. Every bin is checked, including
for empty input and the hand example. Estimate private storage in bytes for
$p$ threads and $b$ bins. Timings include allocation, zeroing, counting and merging;
explain why counting-loop speedup alone does not describe the complete operation.
\end{enumerate}

\exercise{C3}{Recursive tasks and a cutoff}
Edit \code{C3-tasks.cpp}. A range is divided into two parts until its length is
at most \code{cutoff}; short ranges are summed serially. Too many tiny tasks can
cost more to create than they save in arithmetic time.
\begin{enumerate}
\item Draw the split for 11 elements and cutoff 3. Label half-open ranges and
verify that the leaves cover all elements once. An empty range returns zero.
\item Implement \code{split_sum} and \code{task_sum}. Use one outer parallel
region and one \code{single} region to start the root call. Do not create a new
team at each recursive call.
\item For a large range, create one child task for its left half and calculate
the right half in the current task. Copy range parameters with
\code{firstprivate}; share the left result variable with the child that writes it.
Use \code{taskwait} before combining results or returning from the helper.
Explain why the helper's stack must remain alive until the child finishes.
\item Build and run:
\begin{lstlisting}[style=shell]
cmake --build build --target c3_tasks -j
OMP_NUM_THREADS=4 OMP_DYNAMIC=false ./build/bin/c3_tasks 1048576 2048 5
\end{lstlisting}
The arguments are length, positive cutoff and trials. With \code{--self-test}, extra checks include empty,
uneven and short ranges and cutoffs larger than the input. Compare cutoffs
32, 256, 2048 and 16384 at a fixed input and team size. Explain the resulting
balance between available tasks and task-management overhead.
\end{enumerate}

\exercise{C4}{A persistent team across Jacobi steps}
Edit \code{C4-jacobi.cpp}. The owned values are at indices 1 through $n$;
indices 0 and $n+1$ are fixed zeros. Each step computes
\[u_i^{t+1}=\tfrac12(u_{i-1}^{t}+u_{i+1}^{t}).\]
All updates must use the old array. In-place updates would change the algorithm.
\begin{enumerate}
\item From $(0,2,4,8,0)$, derive one step and then a second step.
Draw the two old inputs required by each new value.
\item Implement \code{jacobi} using two buffers and one parallel region around
all steps. Distribute interior writes with a worksharing loop; retain its end
barrier before swapping roles.
\item Swap shared pointers once using \code{single}, retaining its end barrier.
Explain the two distinct ordering requirements: finish writes, then publish the
new pointer roles. Never modify the fixed ends.
\item Return the final result in the caller's vector. Trace zero, one and two
steps; an odd step count can leave the answer in the other buffer.
\item Build and run:
\begin{lstlisting}[style=shell]
cmake --build build --target c4_jacobi -j
OMP_NUM_THREADS=4 OMP_DYNAMIC=false ./build/bin/c4_jacobi 65536 20 5
\end{lstlisting}
The arguments are interior length, steps and trials. With \code{--self-test}, extra checks include no interior
values, one value, and zero, odd and even step counts. The timed operation includes
a fresh copy of the initial vector and your function; validation is excluded.
\end{enumerate}

\exercise{C5}{Optional: false sharing}
Open \code{C5-false-sharing.cpp}. Each thread repeatedly increments its own
atomic counter. Separate counters may still share a coherence line, so race-free
accesses can cause ownership traffic between cores.
\begin{enumerate}
\item Leave \code{Packed} unchanged and add 64-byte alignment to
\code{Separated}. Predict how this changes object spacing for a 64-byte cache-line
model. Check the actual line size before applying the model to the machine.
\item Build and run on distinct allocated cores:
\begin{lstlisting}[style=shell]
cmake --build build --target c5_false_sharing -j
OMP_NUM_THREADS=4 OMP_DYNAMIC=false OMP_PLACES=cores \
  OMP_PROC_BIND=spread ./build/bin/c5_false_sharing 1000000 5
\end{lstlisting}
The arguments are increments per thread and trials. The program checks final
counts and prints every time sample, alignment and object size. Compute the
median for each layout, then compare one thread with several threads.
\item Explain why atomic increments preserve correctness but do not remove
false sharing. Their own cost is included, so this experiment does not measure
ordinary store cost alone. Distinguish a measured effect from a proposed cause.
\end{enumerate}
